\section{Discussion}

The completed campaigns illustrate a central methodological point: stop time, warmup interval, RTT, queue-disc settings, and flow mix jointly shape the interpretation of TCP/AQM results. The earlier smoke run was adequate for checking that the workflow executed, but it under-represented high-RTT behavior. The full campaign therefore makes the stop time and warmup interval explicit and stores them with every run.

The benchmark also makes ECN configuration explicit at both the TCP and queue-disc layers. DCTCP is excluded from the main non-ECN comparison because DCTCP without queue marking is not a meaningful protocol comparison. Such cases can still be useful as warnings or failure-mode checks, but they should not be mixed into the main result matrix.

The full single-flow campaign also reveals a limitation of deterministic traffic. Three random-run values are useful for reproducibility bookkeeping, but they do not create meaningful variability when the scenario does not consume random variables. In the present single-flow campaign, the three \texttt{RngRun} values produce numerically identical traces because no stochastic input is sampled; reported single-flow confidence intervals are therefore structurally zero, not tight, and we choose to disclose this rather than display them. The mixed-flow campaign addresses this by adding a competing flow and \texttt{rngRun}-driven start jitter. The resulting mixed-flow confidence intervals do reflect controlled workload variation. The harness has since been extended with a first-flow \texttt{firstStartJitter} input so that future single-flow campaigns can produce equally meaningful confidence intervals without changing scenario semantics.

The ECN comparison is a useful negative result in the current pilot. For single-flow CUBIC and Reno, enabling ECN changes mean queueing delay by only a small amount under many configurations, while DCTCP depends on ECN marking and is therefore treated separately. This does not mean ECN is unimportant; it means the paper should avoid over-claiming from a single-flow deterministic topology and should emphasize the interaction between TCP semantics, queue marking, and workload mix.

The warmup and late-window audits also change how the final evaluation should be framed. Several 40 s configurations, especially high-RTT DCTCP/PIE and selected mixed-flow cases, do not appear stable enough for final performance claims. Rather than hiding this, the artifact exposes it. The next campaign should therefore be targeted: rerun the unstable configurations with longer stop times and report whether rankings and fairness conclusions survive the longer observation window.

Finally, preserving raw traces changes the reviewability of the experiment. Throughput summaries alone can hide queueing-delay behavior, and aggregate means can hide transient or tail behavior. Keeping raw queue length, mark, drop, RTT, throughput, and congestion-window traces allows later analyses to be added without rerunning the full campaign.
